% Smallest Enclosing Circle (Welzl's Algorithm)

Finds the minimum bounding circle in $O(N)$ expected time. Uses global \lstinline|Point|, \lstinline|dist|, and \lstinline|rng|.

\begin{lstlisting}
struct Circle { Point center; double r; };

bool is_inside(Point p, Circle c) { return dist(p, c.center) <= c.r + 1e-9; }

Circle circle_from_3_points(Point A, Point B, Point C) {
    double bx = B.x - A.x, by = B.y - A.y, cx = C.x - A.x, cy = C.y - A.y;
    double BB = bx * bx + by * by, CC = cx * cx + cy * cy, D = 2 * (bx * cy - by * cx);
    Point center = {A.x + (cy * BB - by * CC) / D, A.y + (bx * CC - cx * BB) / D};
    return {center, dist(center, A)};
}

Circle circle_from_2_points(Point A, Point B) {
    return {{(A.x + B.x) / 2.0, (A.y + B.y) / 2.0}, dist(A, B) / 2.0};
}

Circle welzl(vector<Point>& pts, vector<Point> r, int n) {
    if (n == 0 || r.size() == 3) {
        if (r.empty()) return {{0, 0}, 0};
        if (r.size() == 1) return {r[0], 0};
        if (r.size() == 2) return circle_from_2_points(r[0], r[1]);
        return circle_from_3_points(r[0], r[1], r[2]);
    }
    Point p = pts[n - 1]; Circle c = welzl(pts, r, n - 1);
    if (is_inside(p, c)) return c;
    r.push_back(p); return welzl(pts, r, n - 1);
}

Circle get_smallest_enclosing_circle(vector<Point> pts) {
    shuffle(pts.begin(), pts.end(), rng);
    return welzl(pts, {}, pts.size());
}
\end{lstlisting}
